\setcounter{mtc}{5} %indique le numero reel du chapitre, pour la mini table des matieres
\chapter{etude Theorique}
\minitoc  %insert la minitoc

\graphicspath{{Chapitre1/figures/}}
%==============================================================================
\pagestyle{fancy}
\fancyhf{}
\fancyhead[R]{\bfseries\rightmark}
\fancyfoot[R]{\thepage}
\renewcommand{\headrulewidth}{0.5pt}
\renewcommand{\footrulewidth}{0pt}
\renewcommand{\chaptermark}[1]{\markboth{\MakeUppercase{\chaptername~\thechapter. #1 }}{}}
\renewcommand{\sectionmark}[1]{\markright{\thechapter.\thesection~ #1}}

\begin{spacing}{1.2}
%==============================================================================

\section*{Introduction}
Une etude theorique \cite{YOUSFI2015} peut contenir l'une et/ou l'autre de ces deux parties :
\section{etat de l'art} 
C'est une etude assez detaillee sur ce qui existe sur le marche ou dans la litterature (d'oe 
le terme etat de l'art), qui permet de repondre e la problematique. L'idee ici est de faire 
un comparatif entre les solutions existantes, mais surtout d'analyser le resultat de cette 
comparaison et de dire pourquoi ne sont-elles pas satisfaisantes pour repondre e votre 
problematique.
\begin{figure}[!ht]\centering
\includegraphics[scale=0.9]{art.jpg}
\caption{etat de l'art}
\label{fig:fig1}
\end{figure}
\section{etude de l'existant}
Elle est en general realisee quand on va developper un module supplementaire sur un 
logiciel existant, ou si on va modifier une application existante. L'etude de l'existant
consiste e expliquer ce qui existe deje dans votre environnement de travail.

\section*{Conclusion}
La conclusion est en general sans numerotation, et n'apparaet pas dans la table des matieres.


%==============================================================================
\end{spacing}
